\section{Open Questions}

The following five questions are genuinely unresolved after this replication. Each is
grounded in an observation from the paper re-read, not manufactured for completeness.

\begin{enumerate}

\item \textbf{Does the Fornalski 2022 MC tree, once reconstructed, actually produce the
paper's headline numbers?}\\
\emph{Basis:} Every MC-derived headline (82\% full-repair, up to 7\% two-dose AR,
0.126\% global colony, $\sim\!55\times$ MC erosion of the analytical bound) depends on
Fornalski \emph{et al.} 2022 (\emph{Dose-Response}, DOI 10.1177/15593258221103459)
which has no public code release. The present paper does not fully specify the
per-time-step transition graph between \{healthy, damaged, mutated, cancerous\}
states, nor the death/multiplication/classical-repair probabilities, so it is
impossible to re-implement from \S2 alone. This is the dominant audit blocker.\\
\emph{Next steps:} Read Fornalski 2022 end-to-end plus refs [21,22,24,25]; implement the
per-cell transition graph in NumPy/Numba (64\,000 cells $\times$ 50 MC runs, $T_0=120$~h);
reproduce Fig.~2 single-dose overlay and the 82\% full-repair number as the smallest
self-contained MC test; if within $\pm 10\%$ of the paper, extend to Figs.~3--11 and the
0.126\% global colony contribution. Budget 3--5 days if the parent paper is well
specified, 1--2 weeks otherwise. CPU-only.

\item \textbf{Priming-dose window kinetics: does the analytical
$\exp(-\alpha_2 k)$ decay actually match the MC Fig.~6 ``$\le 80$~h persistence'' claim?}\\
\emph{Basis:} $\alpha_2=0.0832$~h$^{-1}$ implies 24~h $\to$ peak, 48~h $\to$ 13.5\% of
peak, 80~h $\to$ 0.13\% of peak. Fig.~6 is a MC ensemble property that can decay
faster or slower than the analytical envelope depending on cell-division dilution and
death dynamics in the MC tree. Neither confirmed nor refuted by the analytical layer
alone.\\
\emph{Next steps:} Once the MC tree is reconstructed (Q1), sweep $\Delta t$ from 12--240~h
in 12~h steps at fixed $(D_1{=}25\text{~mGy}, D_2{=}1500\text{~mGy})$, record AR-on vs
AR-off survival $\Delta$, fit a decay time constant, compare to
$\alpha_2^{-1}=12.02$~h. If MC decay is substantially slower than the analytical
envelope, that would be a novel finding.

\item \textbf{Is $P_{AR}(D,k)$ reducible to a specific molecular pathway, or is it
mechanism-degenerate?}\\
\emph{Basis:} The paper's ``mechanistic'' modelling is three empirical constants with no
mapping to measurable pathway outputs (foci counts, phospho-ATM, phospho-p53, NRF2
translocation). Any smooth peaked-in-$(D,k)$ function fits the calibration data, so
``mechanism'' claims are underdetermined.\\
\emph{Next steps:} Take a published pathway-level AR model (Schollnberger 2002
stochastic ATM/p53, or a recent NRF2-based model), simulate at a $(D_1,\Delta t)$
grid, fit $(\alpha_0,\alpha_1,\alpha_2)$ to those simulated outputs, check
identifiability: do multiple pathway-level parameter sets give the same
$(\alpha_0,\alpha_1,\alpha_2)$? If yes, Piotrowski \emph{et al.}'s model is degenerate
w.r.t.\ mechanism and the title claim is misleading; if no, a fitted
$(\alpha_0,\alpha_1,\alpha_2)$ is a pathway fingerprint and the paper is doing something
stronger than its text admits.

\item \textbf{Cell-line and LET generalizability of $\alpha_0,\alpha_1,\alpha_2$ and
$\mu_0,\mu_1$.}\\
\emph{Basis:} Paper uses two parameter sets (HBRA in-vivo + in-vitro human lymphocyte
X-ray). No cell-line or LET dependence reported. AR is experimentally known to vary by
order-of-magnitude across cell lines and to weaken sharply above $\sim 1$~keV/$\mu$m
LET. If the constants are cell-line/LET-specific, the model's predictive scope is much
narrower than the paper implies.\\
\emph{Next steps:} Extract published dose-response curves for AR in V79 (Wolff 1996),
CHO (Shadley \& Wolff 1987), TK6 (Zhou 2004), primary human fibroblasts (Sasaki 2006);
re-fit $(\alpha_0,\alpha_1,\alpha_2)$ per cell line; test clustering vs spread. For LET,
use IMRT proton and $^{12}$C-ion AR datasets (Belli 2000). Report whether the paper's
constants are outliers or representative.

\item \textbf{Resolution of the \S3.4 numerical inconsistency (MISMATCH C4c).}\\
\emph{Basis:} Verified in the audit: Eq.~(5) with (in-vitro, 0.17~mGy/h) gives
$P_C=0.155$, and (HBRA, 0.002~mGy/h) gives $P_C=0.454$. Paper attributes
$P_C=0.45$ to scenario~3 (Fig.~9), but only scenario~2's parameters produce it.
Three resolution candidates: (a) scenarios 2--3 swapped in the discussion, (b) bullet-list
parameters wrong for scenario~3, (c) Eq.~(5) coefficients (11.5, 38) are typos. Each
has different implications for whether Figs.~7--10 actually show the right curves.\\
\emph{Next steps:} Pixel-diff Figs.~7 and 9 against re-generated Eq.~(5) curves under
each resolution hypothesis (plot $P_C(\dot{d})$ with in-vitro, HBRA, and each
candidate typo-fix, match visual shape to Fig.~9). Once resolved, email K.W. Fornalski
(corresponding author) with the finding; propose an erratum.

\end{enumerate}
